% 05-planes.tex — Sección 5: planes de ejecución. Responde la pregunta 6 del
% revisor y declara la asimetría de planes (riesgo R14 del PLAN).

\section{Planes de ejecución}
\label{sec:planes}

Se capturó el plan de cada combinación de motor, consulta y perfil con
\texttt{explain} en modo \texttt{executionStats} (MongoDB) y
\texttt{EXPLAIN (ANALYZE, BUFFERS)} (PostgreSQL). Los 12 planes crudos están en
\texttt{results/explains/}.

\subsection{Dos avisos para leerlos}

\paragraph{Los planes se capturan con un solo cliente.}
Se toman en las celdas de \mbox{c=1} y \textbf{fuera} de la ventana de medición,
para no alterar lo que se mide. Eso tiene una consecuencia que hay que declarar:
\textbf{el plan de \mbox{c=1} no es el que se ejecuta con \mbox{c=50}}. Con 50
clientes simultáneos, PostgreSQL agota los procesos de trabajo paralelos
disponibles y los clientes que llegan después ejecutan la consulta sin
paralelismo. El plan documenta la \emph{forma} de resolver la consulta, no el
grado de paralelismo que tuvo cada cliente bajo carga.

\paragraph{En \texttt{EXPLAIN (BUFFERS)} los contadores incluyen a los hijos.}
El total es el del nodo raíz. Sumar los nodos cuenta lo mismo varias veces: al
verificar las consultas se registraron 8310 bloques donde en realidad eran 1662.

\subsection{La asimetría de planes, declarada}
\label{sec:asimetria-planes}

\texttt{Q2} filtra por \texttt{timestamp} y ordena por \texttt{final\_total}. Los
dos motores la resuelven de forma distinta:

\begin{itemize}
  \item \textbf{MongoDB} usa su índice sobre \texttt{timestamp} para el
        \texttt{\$match}.
  \item \textbf{PostgreSQL} hace \texttt{Seq Scan}. El índice de expresión existe
        —se creó sobre el \texttt{timestamp} extraído como texto, sin modificar la
        consulta del manual— pero el planificador no lo usa, porque la mitad de la
        tabla pasa el filtro y recorrerla entera le sale más barato que saltar por
        el índice.
\end{itemize}

Esto es importante al interpretar las latencias: \textbf{en parte se están
comparando dos planes distintos}, no solo dos motores. No se «arregló» forzando
el índice, porque eso significaría medir una consulta que no es la del manual.

\subsection{Lo que el plan explica del cuello de botella (pregunta 6)}

\paragraph{\texttt{Q1} en PostgreSQL: el plan explica el volcado.}
El plan muestra el ordenamiento con \texttt{Sort Method: external merge} y
\texttt{Sort Space Type: Disk}, que es exactamente lo que registran los
contadores \texttt{temp\_bytes} de la sección~\ref{sec:spill}. Las dos fuentes
coinciden: no es una conjetura sobre por qué \texttt{Q1} es lenta en PostgreSQL.

\paragraph{\texttt{Q2} en los dos motores: el plan explica por qué \emph{no} hay volcado.}
El \texttt{LIMIT 1000} convierte el ordenamiento en un \emph{top-K}: no hace falta
ordenar las \num{175049} filas que pasan el filtro, solo mantener las 1000
mejores. PostgreSQL lo resuelve con \texttt{top-N heapsort} y
\texttt{Sort Space Used: 273}~kB \textbf{en memoria}. Esa cifra es la clave de la
sección~\ref{sec:variantes}: 273~kB caben de sobra tanto en los 4~MB de
\texttt{work\_mem} de la variante restrictiva como en los 64~MB de la permisiva,
y por eso las dos variantes de \texttt{Q2} no se diferencian.
